\chapter{Dynamic Programming: 2D/3D DP and DP on Grids}
\chapterepigraph{A 1D DP table has one knob: which earlier position(s) a
state depends on. A grid DP table has two -- row and column -- and the
only genuinely new idea in this chapter is that a cell's recurrence now
reads from \emph{neighbors in two directions} instead of \emph{earlier
points on one line}. Everything else, all four stages, transfers
unchanged.}

\section{From a Line to a Grid}
\label{sec:dp-grids-intro}

Every problem in the DP 1D chapter (Fibonacci through House Robber II)
indexed its DP state by a single integer (a position along an array).
This chapter's problems index state by \textbf{two} integers (a row and
a column), because the underlying problem is itself naturally laid out on
a grid. The four-stage methodology from the DP 1D chapter's introduction --
recursion, memoization, tabulation, space optimization -- applies
identically; only the \emph{shape} of the memo array and the tabulation
loop changes, from a 1D array scanned in one direction to a 2D array
filled row by row (or column by column).

\begin{patternbox}[How to Recognise a Grid DP Problem and Its Recurrence]
\begin{itemize}
  \item \textbf{``Number of paths / minimum-cost path from top-left to
        bottom-right, moving only right or down''} -- the recurrence at
        cell $(i,j)$ combines the answers at $(i-1,j)$ (arrived from
        above) and $(i,j-1)$ (arrived from the left), summed for
        counting paths, or minimized/maximized for cost.
  \item \textbf{``Minimum/maximum path sum from the top of a triangle (or
        any row) to the bottom''} -- the recurrence at $(i,j)$ combines
        the two cells reachable \emph{from} it in the next row (or the
        two reachable \emph{into} it from the previous row, depending on
        which direction the recursion scans), rather than from directly
        above/left -- the exact indices depend on whether you formulate
        top-down or bottom-up first.
  \item \textbf{``Choose one of several options each day/step, but the
        same option cannot be chosen on two consecutive
        days/steps''} -- a state with \emph{two} dimensions where one is
        not spatial at all: the position (day/step) and \emph{which
        option was just used}, so the recurrence can exclude the
        previous day's choice when picking the best current option.
  \item \textbf{``Two agents/paths moving simultaneously through the
        same grid''} -- a \textbf{3D} DP state (one shared row/step index,
        plus one column index per agent), since both agents' positions
        must be tracked together to correctly avoid double-counting
        cells they both visit.
\end{itemize}
\end{patternbox}

This chapter moves through exactly these four variations in order,
starting with the cleanest two-direction grid recurrence (Unique Paths)
and ending with the full 3D, two-agent recurrence (Ninja and his
Friends), which is this chapter's capstone, exactly as Largest BST in a
Binary Tree and the LFU Cache closed their respective chapters.

\newpage

\section{Unique Paths}
\label{sec:dp-unique-paths}

\problemstatement{A robot starts at the top-left corner of an
$m \times n$ grid and wants to reach the bottom-right corner, moving only
\textbf{right} or \textbf{down} at each step. Count the number of unique
paths.}

\begin{patternbox}
\emph{``Number of paths through a grid, moving only right or down''} is
the cleanest possible grid-DP recurrence: a path arrives at cell
$(i,j)$ from exactly one of two neighbors -- directly above, $(i-1,j)$,
or directly to the left, $(i,j-1)$ -- and since every path to $(i,j)$ is
either one of these two cases, the count of paths to $(i,j)$ is their
\textbf{sum}, exactly mirroring the DP 1D chapter's Climbing Stairs
``sum over a fixed set of predecessors,'' just with the predecessors now
arranged on two grid axes instead of one line.
\end{patternbox}

\subsection*{Deriving the recurrence}

Let $\textit{ways}(i,j)$ be the number of unique paths from $(0,0)$ to
$(i,j)$. Since the only moves are right and down, the last move into
$(i,j)$ came from $(i-1,j)$ (a down-move) or $(i,j-1)$ (a right-move):
\[
  \textit{ways}(i,j) = \textit{ways}(i-1,j) + \textit{ways}(i,j-1).
\]
Base case: $\textit{ways}(0,0) = 1$ (the trivial path of staying put);
any $\textit{ways}(i,j)$ with $i<0$ or $j<0$ is treated as $0$
(non-existent, out-of-grid predecessors contribute nothing).

\subsection*{Approach 1: Recursion}

\begin{lstlisting}
#include <bits/stdc++.h>
using namespace std;

int uniquePathsRecursive(int i, int j) {
    if (i == 0 && j == 0) return 1;
    if (i < 0 || j < 0) return 0;

    int fromAbove = uniquePathsRecursive(i - 1, j);
    int fromLeft  = uniquePathsRecursive(i, j - 1);
    return fromAbove + fromLeft;
}

int main() {
    int m = 3, n = 7;
    cout << uniquePathsRecursive(m - 1, n - 1) << endl; // 28
    return 0;
}
\end{lstlisting}

\begin{complexitybox}[Approach 1 Complexity]
\textbf{Time.} $O(2^{m+n})$ -- each call branches into two, with
recursion depth up to $m+n$.
\textbf{Space.} $O(m+n)$ for the recursion stack.
\end{complexitybox}

\subsection*{Approach 2: Memoization}

\begin{lstlisting}
#include <bits/stdc++.h>
using namespace std;

int uniquePathsMemo(int i, int j, vector<vector<int>>& memo) {
    if (i == 0 && j == 0) return 1;
    if (i < 0 || j < 0) return 0;
    if (memo[i][j] != -1) return memo[i][j];

    int fromAbove = uniquePathsMemo(i - 1, j, memo);
    int fromLeft  = uniquePathsMemo(i, j - 1, memo);
    return memo[i][j] = fromAbove + fromLeft;
}

int main() {
    int m = 3, n = 7;
    vector<vector<int>> memo(m, vector<int>(n, -1));
    cout << uniquePathsMemo(m - 1, n - 1, memo) << endl; // 28
    return 0;
}
\end{lstlisting}

\begin{complexitybox}[Approach 2 Complexity]
\textbf{Time.} $O(m \cdot n)$ -- one $O(1)$ computation per distinct
$(i,j)$ pair.
\textbf{Space.} $O(m \cdot n)$ memo $+$ $O(m+n)$ recursion stack.
\end{complexitybox}

\subsection*{Approach 3: Tabulation}

\begin{lstlisting}
#include <bits/stdc++.h>
using namespace std;

int uniquePathsTabulation(int m, int n) {
    vector<vector<int>> dp(m, vector<int>(n, 0));

    for (int i = 0; i < m; i++) {
        for (int j = 0; j < n; j++) {
            if (i == 0 && j == 0) {
                dp[i][j] = 1;
                continue;
            }
            int fromAbove = (i > 0) ? dp[i-1][j] : 0;
            int fromLeft  = (j > 0) ? dp[i][j-1] : 0;
            dp[i][j] = fromAbove + fromLeft;
        }
    }
    return dp[m - 1][n - 1];
}

int main() {
    cout << uniquePathsTabulation(3, 7) << endl; // 28
    return 0;
}
\end{lstlisting}

\subsection*{Dry run of the tabulation table}

\dryrun{Take $m=3$, $n=3$.}

\begin{center}
\begin{tabular}{c|c|c|c}
 & $j=0$ & $j=1$ & $j=2$ \\ \hline
$i=0$ & 1 & 1 & 1 \\ \hline
$i=1$ & 1 & 2 & 3 \\ \hline
$i=2$ & 1 & 3 & 6
\end{tabular}
\end{center}

The first row and first column are all $1$s (only one way to reach any
cell along the top edge or left edge: a straight line of moves). Every
other cell is the sum of its top and left neighbors, e.g.\ $dp[2][2] =
dp[1][2]+dp[2][1] = 3+3=6$. The answer for a $3\times3$ grid is $6$.

\begin{complexitybox}[Approach 3 Complexity]
\textbf{Time.} $O(m \cdot n)$.
\textbf{Space.} $O(m \cdot n)$ table, no recursion stack.
\end{complexitybox}

\subsection*{Approach 4: Space Optimization}

\begin{ideabox}[Only the Previous Row Is Ever Needed]
Computing row $i$ only reads from row $i-1$ (directly above) and from
\emph{within} row $i$ itself (directly left, already computed earlier in
the same row's left-to-right scan). No row before $i-1$ is ever needed
again once row $i$ is being filled, so a single 1D array, reused and
updated in place across rows, suffices.
\end{ideabox}

\begin{lstlisting}
#include <bits/stdc++.h>
using namespace std;

int uniquePathsSpaceOptimized(int m, int n) {
    vector<int> prevRow(n, 0);

    for (int i = 0; i < m; i++) {
        vector<int> currentRow(n, 0);
        for (int j = 0; j < n; j++) {
            if (i == 0 && j == 0) {
                currentRow[j] = 1;
                continue;
            }
            int fromAbove = prevRow[j];
            int fromLeft  = (j > 0) ? currentRow[j-1] : 0;
            currentRow[j] = fromAbove + fromLeft;
        }
        prevRow = currentRow;
    }
    return prevRow[n - 1];
}

int main() {
    cout << uniquePathsSpaceOptimized(3, 7) << endl; // 28
    return 0;
}
\end{lstlisting}

\begin{complexitybox}[Approach 4 Complexity]
\textbf{Time.} $O(m \cdot n)$.
\textbf{Space.} $O(n)$ for the single rolling row, instead of
$O(m \cdot n)$ for the full table.
\end{complexitybox}

\begin{takeawaybox}
Grid DP's space optimization mirrors 1D DP's exactly, just one dimension
up: instead of two rolling \emph{scalars} (Fibonacci, House Robber), we
keep two rolling \emph{rows} (or, with more care, overwrite a single row
in place, since the left-neighbor read within a row happens before that
position is overwritten in a left-to-right scan). Whenever a grid
recurrence only reads from the immediately preceding row and the current
row, this reduction from $O(mn)$ to $O(n)$ space is always available.
\end{takeawaybox}

\section{Unique Paths II (With Obstacles)}
\label{sec:dp-unique-paths-2}

\problemstatement{Same as Section~\ref{sec:dp-unique-paths}, except some
grid cells are marked as \textbf{obstacles} (value $1$; open cells are
$0$), and the robot cannot step onto an obstacle cell at all. Count the
number of unique paths from top-left to bottom-right.}

\begin{patternbox}
\emph{``Same grid-path-counting problem, with some cells forbidden''}
needs exactly one change to Section~\ref{sec:dp-unique-paths}'s
recurrence: an obstacle cell contributes \textbf{zero} paths, regardless
of what its predecessors computed, because no path is allowed to pass
through it at all.
\end{patternbox}

\subsection*{Deriving the recurrence}

\[
  \textit{ways}(i,j) =
  \begin{cases}
    0 & \text{if } (i,j) \text{ is an obstacle}, \\
    \textit{ways}(i-1,j) + \textit{ways}(i,j-1) & \text{otherwise.}
  \end{cases}
\]
The base case also needs the obstacle check: $\textit{ways}(0,0) = 1$
\emph{unless} $(0,0)$ itself is an obstacle, in which case no path can
even begin, and the answer is $0$ overall.

\subsection*{Approach 1: Recursion}

\begin{lstlisting}
#include <bits/stdc++.h>
using namespace std;

int uniquePaths2Recursive(int i, int j, vector<vector<int>>& grid) {
    if (i < 0 || j < 0 || grid[i][j] == 1) return 0;
    if (i == 0 && j == 0) return 1;

    int fromAbove = uniquePaths2Recursive(i - 1, j, grid);
    int fromLeft  = uniquePaths2Recursive(i, j - 1, grid);
    return fromAbove + fromLeft;
}

int main() {
    vector<vector<int>> grid = {{0,0,0},{0,1,0},{0,0,0}};
    int m = grid.size(), n = grid[0].size();
    cout << uniquePaths2Recursive(m - 1, n - 1, grid) << endl; // 2
    return 0;
}
\end{lstlisting}

\begin{complexitybox}[Approach 1 Complexity]
\textbf{Time.} $O(2^{m+n})$.
\textbf{Space.} $O(m+n)$ recursion stack.
\end{complexitybox}

\subsection*{Approach 2: Memoization}

\begin{lstlisting}
#include <bits/stdc++.h>
using namespace std;

int uniquePaths2Memo(int i, int j, vector<vector<int>>& grid,
                      vector<vector<int>>& memo) {
    if (i < 0 || j < 0 || grid[i][j] == 1) return 0;
    if (i == 0 && j == 0) return 1;
    if (memo[i][j] != -1) return memo[i][j];

    int fromAbove = uniquePaths2Memo(i - 1, j, grid, memo);
    int fromLeft  = uniquePaths2Memo(i, j - 1, grid, memo);
    return memo[i][j] = fromAbove + fromLeft;
}

int main() {
    vector<vector<int>> grid = {{0,0,0},{0,1,0},{0,0,0}};
    int m = grid.size(), n = grid[0].size();
    vector<vector<int>> memo(m, vector<int>(n, -1));
    cout << uniquePaths2Memo(m - 1, n - 1, grid, memo) << endl; // 2
    return 0;
}
\end{lstlisting}

\begin{complexitybox}[Approach 2 Complexity]
\textbf{Time.} $O(m \cdot n)$.
\textbf{Space.} $O(m \cdot n)$ memo $+$ $O(m+n)$ recursion stack.
\end{complexitybox}

\subsection*{Approach 3: Tabulation}

\begin{lstlisting}
#include <bits/stdc++.h>
using namespace std;

int uniquePaths2Tabulation(vector<vector<int>>& grid) {
    int m = grid.size(), n = grid[0].size();
    vector<vector<int>> dp(m, vector<int>(n, 0));

    for (int i = 0; i < m; i++) {
        for (int j = 0; j < n; j++) {
            if (grid[i][j] == 1) {
                dp[i][j] = 0;
                continue;
            }
            if (i == 0 && j == 0) {
                dp[i][j] = 1;
                continue;
            }
            int fromAbove = (i > 0) ? dp[i-1][j] : 0;
            int fromLeft  = (j > 0) ? dp[i][j-1] : 0;
            dp[i][j] = fromAbove + fromLeft;
        }
    }
    return dp[m - 1][n - 1];
}

int main() {
    vector<vector<int>> grid = {{0,0,0},{0,1,0},{0,0,0}};
    cout << uniquePaths2Tabulation(grid) << endl; // 2
    return 0;
}
\end{lstlisting}

\subsection*{Dry run of the tabulation table}

\dryrun{Take $\textit{grid} = \begin{bmatrix}0&0&0\\0&1&0\\0&0&0\end{bmatrix}$.}

\begin{center}
\begin{tabular}{c|c|c|c}
 & $j=0$ & $j=1$ & $j=2$ \\ \hline
$i=0$ & 1 & 1 & 1 \\ \hline
$i=1$ & 1 & 0 (obstacle) & 1 \\ \hline
$i=2$ & 1 & 1 & 2
\end{tabular}
\end{center}

At $(1,1)$, the obstacle forces $\textit{dp}=0$ regardless of its
neighbors. At $(2,2)$: $\textit{dp}[1][2]+\textit{dp}[2][1] = 1+1=2$.
The answer is $2$: the two paths route either entirely along the top
then down the right edge, or entirely down the left edge then along the
bottom.

\begin{complexitybox}[Approach 3 Complexity]
\textbf{Time.} $O(m \cdot n)$.
\textbf{Space.} $O(m \cdot n)$ table, no recursion stack.
\end{complexitybox}

\subsection*{Approach 4: Space Optimization}

\begin{lstlisting}
#include <bits/stdc++.h>
using namespace std;

int uniquePaths2SpaceOptimized(vector<vector<int>>& grid) {
    int m = grid.size(), n = grid[0].size();
    vector<int> prevRow(n, 0);

    for (int i = 0; i < m; i++) {
        vector<int> currentRow(n, 0);
        for (int j = 0; j < n; j++) {
            if (grid[i][j] == 1) {
                currentRow[j] = 0;
                continue;
            }
            if (i == 0 && j == 0) {
                currentRow[j] = 1;
                continue;
            }
            int fromAbove = prevRow[j];
            int fromLeft  = (j > 0) ? currentRow[j-1] : 0;
            currentRow[j] = fromAbove + fromLeft;
        }
        prevRow = currentRow;
    }
    return prevRow[n - 1];
}

int main() {
    vector<vector<int>> grid = {{0,0,0},{0,1,0},{0,0,0}};
    cout << uniquePaths2SpaceOptimized(grid) << endl; // 2
    return 0;
}
\end{lstlisting}

\begin{complexitybox}[Approach 4 Complexity]
\textbf{Time.} $O(m \cdot n)$.
\textbf{Space.} $O(n)$ for the rolling row.
\end{complexitybox}

\begin{takeawaybox}
Adding a ``forbidden state'' constraint to an already-derived DP
recurrence is almost always a single extra case at the very top of the
recurrence (here, ``obstacle $\Rightarrow$ zero,'' checked before
anything else), with every other part of the four-stage derivation
unchanged. This is the same kind of small, composable extension already
seen when Combination Sum III added a length constraint to Combination
Sum II's base case.
\end{takeawaybox}

\section{Minimum Path Sum}
\label{sec:dp-minimum-path-sum}

\problemstatement{Given an $m \times n$ grid of non-negative numbers,
find a path from the top-left to the bottom-right corner, moving only
right or down, that \textbf{minimizes} the sum of the numbers along the
path.}

\begin{patternbox}
This is Unique Paths (Section~\ref{sec:dp-unique-paths}) with sum
replaced by minimum -- exactly the same transformation that turned
Climbing Stairs into Frog Jump back in the DP 1D chapter. The
recurrence's \emph{structure} (combine the cell above and the cell to
the left) is identical; only the combination operator and the addition
of the current cell's own value change.
\end{patternbox}

\subsection*{Deriving the recurrence}

Let $\textit{minSum}(i,j)$ be the minimum path sum to reach $(i,j)$. The
path's last step into $(i,j)$ came from $(i-1,j)$ or $(i,j-1)$, and we
must add the current cell's own value to whichever predecessor gives the
smaller total:
\[
  \textit{minSum}(i,j) = \textit{grid}[i][j] + \min\big(
  \textit{minSum}(i-1,j),\ \textit{minSum}(i,j-1)\big).
\]
Base case: $\textit{minSum}(0,0) = \textit{grid}[0][0]$. For cells along
the top row or left column, only one predecessor exists (the other is
out of bounds and treated as $+\infty$, never the chosen minimum).

\subsection*{Approach 1: Recursion}

\begin{lstlisting}
#include <bits/stdc++.h>
using namespace std;

int minPathSumRecursive(int i, int j, vector<vector<int>>& grid) {
    if (i == 0 && j == 0) return grid[0][0];
    if (i < 0 || j < 0) return INT_MAX;

    int fromAbove = minPathSumRecursive(i - 1, j, grid);
    int fromLeft  = minPathSumRecursive(i, j - 1, grid);
    return grid[i][j] + min(fromAbove, fromLeft);
}

int main() {
    vector<vector<int>> grid = {{1,3,1},{1,5,1},{4,2,1}};
    int m = grid.size(), n = grid[0].size();
    cout << minPathSumRecursive(m - 1, n - 1, grid) << endl; // 7
    return 0;
}
\end{lstlisting}

\begin{complexitybox}[Approach 1 Complexity]
\textbf{Time.} $O(2^{m+n})$.
\textbf{Space.} $O(m+n)$ recursion stack.
\end{complexitybox}

\subsection*{Approach 2: Memoization}

\begin{lstlisting}
#include <bits/stdc++.h>
using namespace std;

int minPathSumMemo(int i, int j, vector<vector<int>>& grid,
                    vector<vector<int>>& memo) {
    if (i == 0 && j == 0) return grid[0][0];
    if (i < 0 || j < 0) return INT_MAX;
    if (memo[i][j] != -1) return memo[i][j];

    int fromAbove = minPathSumMemo(i - 1, j, grid, memo);
    int fromLeft  = minPathSumMemo(i, j - 1, grid, memo);
    return memo[i][j] = grid[i][j] + min(fromAbove, fromLeft);
}

int main() {
    vector<vector<int>> grid = {{1,3,1},{1,5,1},{4,2,1}};
    int m = grid.size(), n = grid[0].size();
    vector<vector<int>> memo(m, vector<int>(n, -1));
    cout << minPathSumMemo(m - 1, n - 1, grid, memo) << endl; // 7
    return 0;
}
\end{lstlisting}

\begin{complexitybox}[Approach 2 Complexity]
\textbf{Time.} $O(m \cdot n)$.
\textbf{Space.} $O(m \cdot n)$ memo $+$ $O(m+n)$ recursion stack.
\end{complexitybox}

\subsection*{Approach 3: Tabulation}

\begin{lstlisting}
#include <bits/stdc++.h>
using namespace std;

int minPathSumTabulation(vector<vector<int>>& grid) {
    int m = grid.size(), n = grid[0].size();
    vector<vector<int>> dp(m, vector<int>(n, 0));

    for (int i = 0; i < m; i++) {
        for (int j = 0; j < n; j++) {
            if (i == 0 && j == 0) {
                dp[i][j] = grid[0][0];
                continue;
            }
            int fromAbove = (i > 0) ? dp[i-1][j] : INT_MAX;
            int fromLeft  = (j > 0) ? dp[i][j-1] : INT_MAX;
            dp[i][j] = grid[i][j] + min(fromAbove, fromLeft);
        }
    }
    return dp[m - 1][n - 1];
}

int main() {
    vector<vector<int>> grid = {{1,3,1},{1,5,1},{4,2,1}};
    cout << minPathSumTabulation(grid) << endl; // 7
    return 0;
}
\end{lstlisting}

\subsection*{Dry run of the tabulation table}

\dryrun{Take $\textit{grid} = \begin{bmatrix}1&3&1\\1&5&1\\4&2&1\end{bmatrix}$.}

\begin{center}
\begin{tabular}{c|c|c|c}
 & $j=0$ & $j=1$ & $j=2$ \\ \hline
$i=0$ & 1 & 4 & 5 \\ \hline
$i=1$ & 2 & 7 & 6 \\ \hline
$i=2$ & 6 & 8 & 7
\end{tabular}
\end{center}

E.g.\ $\textit{dp}[1][2] = \textit{grid}[1][2] + \min(\textit{dp}[0][2],
\textit{dp}[1][1]) = 1+\min(5,7) = 1+5=6$. The answer is
$\textit{dp}[2][2]=7$, achieved by the path $1\to3\to1\to1\to1$ (right,
right, down, down), summing to $1+3+1+1+1=7$.

\begin{complexitybox}[Approach 3 Complexity]
\textbf{Time.} $O(m \cdot n)$.
\textbf{Space.} $O(m \cdot n)$ table, no recursion stack.
\end{complexitybox}

\subsection*{Approach 4: Space Optimization}

\begin{lstlisting}
#include <bits/stdc++.h>
using namespace std;

int minPathSumSpaceOptimized(vector<vector<int>>& grid) {
    int m = grid.size(), n = grid[0].size();
    vector<int> prevRow(n, 0);

    for (int i = 0; i < m; i++) {
        vector<int> currentRow(n, 0);
        for (int j = 0; j < n; j++) {
            if (i == 0 && j == 0) {
                currentRow[j] = grid[0][0];
                continue;
            }
            int fromAbove = prevRow[j];
            int fromLeft  = (j > 0) ? currentRow[j-1] : INT_MAX;
            currentRow[j] = grid[i][j] + min(fromAbove, fromLeft);
        }
        prevRow = currentRow;
    }
    return prevRow[n - 1];
}

int main() {
    vector<vector<int>> grid = {{1,3,1},{1,5,1},{4,2,1}};
    cout << minPathSumSpaceOptimized(grid) << endl; // 7
    return 0;
}
\end{lstlisting}

\begin{complexitybox}[Approach 4 Complexity]
\textbf{Time.} $O(m \cdot n)$.
\textbf{Space.} $O(n)$ for the rolling row.
\end{complexitybox}

\begin{takeawaybox}
Once a grid-path-counting recurrence is in hand, converting it into a
grid-path-\emph{cost} recurrence is purely mechanical: replace the sum
of predecessors with a minimum (or maximum) of predecessors, plus the
current cell's own contribution. Every later stage -- memoization,
tabulation, space optimization -- carries over with zero additional
changes, exactly as it did for the 1D Fibonacci-to-Frog-Jump
transformation earlier in this book.
\end{takeawaybox}

\section{Triangle}
\label{sec:dp-triangle}

\problemstatement{Given a triangular array of numbers (row $i$ has
$i+1$ entries, $0$-indexed), find the minimum path sum from the top to
any entry of the bottom row, where moving from $(i,j)$ to the next row
may go to $(i+1,j)$ or $(i+1,j+1)$.}

\begin{patternbox}
\emph{``Path through a triangular grid, moving to one of two cells in
the next row''} is structurally the same combine-two-predecessors
recurrence as Minimum Path Sum (Section~\ref{sec:dp-minimum-path-sum}),
but with a twist worth noticing carefully: each cell's choices point
\textbf{forward and downward} (to the next row), not backward
(from the previous row), which means the most natural recursion starts
at the \textbf{top} and recurses toward the bottom, and -- correspondingly
-- the tabulation table is most naturally filled starting from the
\textbf{bottom} row upward, the opposite direction from every grid DP
problem so far in this chapter.
\end{patternbox}

\subsection*{Deriving the recurrence}

Let $\textit{minSum}(i,j)$ be the minimum path sum starting from cell
$(i,j)$ down to \emph{any} cell in the last row. From $(i,j)$, the path
must continue to $(i+1,j)$ or $(i+1,j+1)$, and we want the cheaper
continuation:
\[
  \textit{minSum}(i,j) = \textit{triangle}[i][j] + \min\big(
  \textit{minSum}(i+1,j),\ \textit{minSum}(i+1,j+1)\big).
\]
Base case: at the last row ($i = n-1$), there is nowhere further to go,
so $\textit{minSum}(n-1,j) = \textit{triangle}[n-1][j]$ directly. The
overall answer is $\textit{minSum}(0,0)$.

\subsection*{Approach 1: Recursion}

\begin{lstlisting}
#include <bits/stdc++.h>
using namespace std;

int triangleRecursive(int i, int j, vector<vector<int>>& triangle) {
    int n = triangle.size();
    if (i == n - 1) return triangle[i][j];

    int down      = triangleRecursive(i + 1, j, triangle);
    int downRight = triangleRecursive(i + 1, j + 1, triangle);
    return triangle[i][j] + min(down, downRight);
}

int main() {
    vector<vector<int>> triangle = {{2},{3,4},{6,5,7},{4,1,8,3}};
    cout << triangleRecursive(0, 0, triangle) << endl; // 11
    return 0;
}
\end{lstlisting}

\begin{complexitybox}[Approach 1 Complexity]
\textbf{Time.} $O(2^n)$ -- two branches per call, depth $n$.
\textbf{Space.} $O(n)$ recursion stack.
\end{complexitybox}

\subsection*{Approach 2: Memoization}

\begin{lstlisting}
#include <bits/stdc++.h>
using namespace std;

int triangleMemo(int i, int j, vector<vector<int>>& triangle,
                  vector<vector<int>>& memo) {
    int n = triangle.size();
    if (i == n - 1) return triangle[i][j];
    if (memo[i][j] != -1) return memo[i][j];

    int down      = triangleMemo(i + 1, j, triangle, memo);
    int downRight = triangleMemo(i + 1, j + 1, triangle, memo);
    return memo[i][j] = triangle[i][j] + min(down, downRight);
}

int main() {
    vector<vector<int>> triangle = {{2},{3,4},{6,5,7},{4,1,8,3}};
    int n = triangle.size();
    vector<vector<int>> memo(n, vector<int>(n, -1));
    cout << triangleMemo(0, 0, triangle, memo) << endl; // 11
    return 0;
}
\end{lstlisting}

\begin{complexitybox}[Approach 2 Complexity]
\textbf{Time.} $O(n^2)$ -- the triangle has $O(n^2)$ total cells.
\textbf{Space.} $O(n^2)$ memo $+$ $O(n)$ recursion stack.
\end{complexitybox}

\subsection*{Approach 3: Tabulation}

\begin{ideabox}[Fill the Table Bottom-Up, in the Direction of the Base Case]
Because the recurrence reads from row $i+1$ to compute row $i$, the
tabulation loop must process rows in \textbf{decreasing} order of $i$,
starting by copying the last row directly (the base case), and working
upward to row $0$ -- the mirror image of Sections~\ref{sec:dp-unique-paths}--\ref{sec:dp-minimum-path-sum}'s
top-to-bottom fill order.
\end{ideabox}

\begin{lstlisting}
#include <bits/stdc++.h>
using namespace std;

int triangleTabulation(vector<vector<int>>& triangle) {
    int n = triangle.size();
    vector<vector<int>> dp(n, vector<int>(n, 0));

    for (int j = 0; j < n; j++) dp[n-1][j] = triangle[n-1][j];

    for (int i = n - 2; i >= 0; i--) {
        for (int j = 0; j <= i; j++) {
            int down      = dp[i+1][j];
            int downRight = dp[i+1][j+1];
            dp[i][j] = triangle[i][j] + min(down, downRight);
        }
    }
    return dp[0][0];
}

int main() {
    vector<vector<int>> triangle = {{2},{3,4},{6,5,7},{4,1,8,3}};
    cout << triangleTabulation(triangle) << endl; // 11
    return 0;
}
\end{lstlisting}

\subsection*{Dry run of the tabulation table}

\dryrun{Take $\textit{triangle} = [[2],[3,4],[6,5,7],[4,1,8,3]]$.}

\begin{itemize}
  \item Row $3$ (base case): $\textit{dp}[3] = [4,1,8,3]$.
  \item Row $2$: $\textit{dp}[2][0] = 6+\min(4,1)=7$;
        $\textit{dp}[2][1] = 5+\min(1,8)=6$; $\textit{dp}[2][2] =
        7+\min(8,3)=10$.
  \item Row $1$: $\textit{dp}[1][0] = 3+\min(7,6)=9$;
        $\textit{dp}[1][1] = 4+\min(6,10)=10$.
  \item Row $0$: $\textit{dp}[0][0] = 2+\min(9,10)=11$.
\end{itemize}

The minimum path sum is $\boxed{11}$, via $2\to3\to5\to1$.

\begin{complexitybox}[Approach 3 Complexity]
\textbf{Time.} $O(n^2)$.
\textbf{Space.} $O(n^2)$ table, no recursion stack.
\end{complexitybox}

\subsection*{Approach 4: Space Optimization}

\begin{lstlisting}
#include <bits/stdc++.h>
using namespace std;

int triangleSpaceOptimized(vector<vector<int>>& triangle) {
    int n = triangle.size();
    vector<int> nextRow = triangle[n - 1];

    for (int i = n - 2; i >= 0; i--) {
        vector<int> currentRow(i + 1);
        for (int j = 0; j <= i; j++) {
            int down      = nextRow[j];
            int downRight = nextRow[j + 1];
            currentRow[j] = triangle[i][j] + min(down, downRight);
        }
        nextRow = currentRow;
    }
    return nextRow[0];
}

int main() {
    vector<vector<int>> triangle = {{2},{3,4},{6,5,7},{4,1,8,3}};
    cout << triangleSpaceOptimized(triangle) << endl; // 11
    return 0;
}
\end{lstlisting}

\begin{complexitybox}[Approach 4 Complexity]
\textbf{Time.} $O(n^2)$.
\textbf{Space.} $O(n)$ for the single rolling row.
\end{complexitybox}

\begin{takeawaybox}
Always let the \emph{base case's location} dictate the tabulation's fill
direction -- when a recurrence looks ``forward'' to a later row (as
here), the table must be filled starting from that later base case and
working backward, not from index $0$ forward. Checking which direction
the recursion's base case sits in, before writing the tabulation loop, is
a habit worth forming explicitly to avoid reading uninitialized table
entries.
\end{takeawaybox}

\section{Ninja's Training}
\label{sec:dp-ninja-training}

\problemstatement{A ninja has $n$ days of training. On each day, they
must perform exactly one of $3$ activities (indexed $0,1,2$), earning
$\textit{points}[\textit{day}][\textit{activity}]$ points -- but they may
\textbf{not} perform the \textbf{same} activity on two consecutive days.
Maximize the total points earned over all $n$ days.}

\begin{patternbox}
\emph{``Choose one of several options each day, but not the same option
as the previous day''} is the keyword pattern from
Section~\ref{sec:dp-grids-intro} for a DP state with a \textbf{non-spatial
second dimension}: the state needs both \emph{which day} we are on and
\emph{which activity was excluded} (because it was used the day before),
even though there is no grid or spatial structure here at all -- the
``two dimensions'' are day and constraint, not row and column.
\end{patternbox}

\subsection*{Deriving the recurrence}

Let $f(\textit{day}, \textit{last})$ be the maximum points earnable from
day $0$ through day $\textit{day}$, given that activity $\textit{last}$
is forbidden on day $\textit{day}$ specifically (because it was used on
day $\textit{day}+1$, which we are building backward from -- or, viewed
forward, because it was the most recently committed activity choice we
must avoid repeating). For each candidate activity $\textit{task} \ne
\textit{last}$:
\[
  f(\textit{day}, \textit{last}) = \max_{\textit{task} \ne \textit{last}}
  \Big(\textit{points}[\textit{day}][\textit{task}] +
  f(\textit{day}-1, \textit{task})\Big).
\]
Base case: $f(0, \textit{last}) = \max_{\textit{task} \ne \textit{last}}
\textit{points}[0][\textit{task}]$ (no earlier day to recurse into). The
final answer is $f(n-1, 3)$, using a sentinel value $3$ (outside the
valid activity range $\{0,1,2\}$) to mean ``no activity is forbidden on
the last day,'' since nothing comes after it.

\subsection*{Approach 1: Recursion}

\begin{lstlisting}
#include <bits/stdc++.h>
using namespace std;

int ninjaTrainingRecursive(int day, int last, vector<vector<int>>& points) {
    if (day == 0) {
        int best = 0;
        for (int task = 0; task < 3; task++) {
            if (task != last) best = max(best, points[0][task]);
        }
        return best;
    }

    int best = 0;
    for (int task = 0; task < 3; task++) {
        if (task != last) {
            int candidate = points[day][task] +
                ninjaTrainingRecursive(day - 1, task, points);
            best = max(best, candidate);
        }
    }
    return best;
}

int main() {
    vector<vector<int>> points = {{10,40,70},{20,50,80},{30,60,90}};
    int n = points.size();
    cout << ninjaTrainingRecursive(n - 1, 3, points) << endl; // 210
    return 0;
}
\end{lstlisting}

\begin{complexitybox}[Approach 1 Complexity]
\textbf{Time.} $O(3^n)$ -- up to $3$ branches per call, depth $n$.
\textbf{Space.} $O(n)$ recursion stack.
\end{complexitybox}

\subsection*{Approach 2: Memoization}

\begin{lstlisting}
#include <bits/stdc++.h>
using namespace std;

int ninjaTrainingMemo(int day, int last, vector<vector<int>>& points,
                       vector<vector<int>>& memo) {
    if (day == 0) {
        int best = 0;
        for (int task = 0; task < 3; task++) {
            if (task != last) best = max(best, points[0][task]);
        }
        return best;
    }
    if (memo[day][last] != -1) return memo[day][last];

    int best = 0;
    for (int task = 0; task < 3; task++) {
        if (task != last) {
            int candidate = points[day][task] +
                ninjaTrainingMemo(day - 1, task, points, memo);
            best = max(best, candidate);
        }
    }
    return memo[day][last] = best;
}

int main() {
    vector<vector<int>> points = {{10,40,70},{20,50,80},{30,60,90}};
    int n = points.size();
    vector<vector<int>> memo(n, vector<int>(4, -1));
    cout << ninjaTrainingMemo(n - 1, 3, points, memo) << endl; // 210
    return 0;
}
\end{lstlisting}

\begin{complexitybox}[Approach 2 Complexity]
\textbf{Time.} $O(n \cdot 4 \cdot 3) = O(n)$ -- $n$ days, $4$ possible
\textit{last} values, $3$ choices of \textit{task} per state.
\textbf{Space.} $O(n)$ memo (the $4$ factor is a constant) $+$ $O(n)$
recursion stack.
\end{complexitybox}

\subsection*{Approach 3: Tabulation}

\begin{lstlisting}
#include <bits/stdc++.h>
using namespace std;

int ninjaTrainingTabulation(vector<vector<int>>& points) {
    int n = points.size();
    vector<vector<int>> dp(n, vector<int>(4, 0));

    for (int last = 0; last < 4; last++) {
        int best = 0;
        for (int task = 0; task < 3; task++) {
            if (task != last) best = max(best, points[0][task]);
        }
        dp[0][last] = best;
    }

    for (int day = 1; day < n; day++) {
        for (int last = 0; last < 4; last++) {
            int best = 0;
            for (int task = 0; task < 3; task++) {
                if (task != last) {
                    int candidate = points[day][task] + dp[day-1][task];
                    best = max(best, candidate);
                }
            }
            dp[day][last] = best;
        }
    }
    return dp[n - 1][3];
}

int main() {
    vector<vector<int>> points = {{10,40,70},{20,50,80},{30,60,90}};
    cout << ninjaTrainingTabulation(points) << endl; // 210
    return 0;
}
\end{lstlisting}

\subsection*{Dry run of the tabulation table}

\dryrun{Take $\textit{points} = [[10,40,70],[20,50,80],[30,60,90]]$.}

\begin{center}
\begin{tabular}{c|c|c|c|c}
$\textit{day}$ & $\textit{last}=0$ & $\textit{last}=1$ &
$\textit{last}=2$ & $\textit{last}=3$ \\ \hline
$0$ & 70 & 70 & 40 & 70 \\ \hline
$1$ & 120 & 120 & 120 & 120 \\ \hline
$2$ & 210 & 210 & 180 & 210
\end{tabular}
\end{center}

Reading $\textit{dp}[0][2]=40$: with activity $2$ forbidden on day $0$,
the best of $\{\textit{points}[0][0],\textit{points}[0][1]\} =
\{10,40\}$ is $40$. Reading $\textit{dp}[1][3]$ (no restriction on day
$1$): $\max(20+\textit{dp}[0][0],\ 50+\textit{dp}[0][1],\
80+\textit{dp}[0][2]) = \max(20{+}70,\ 50{+}70,\ 80{+}40) =
\max(90,120,120) = 120$. Reading the final answer,
$\textit{dp}[2][3]$: $\max(30+\textit{dp}[1][0],\ 60+\textit{dp}[1][1],\
90+\textit{dp}[1][2]) = \max(30{+}120,\ 60{+}120,\ 90{+}120) =
\max(150,180,210) = 210$.

\begin{complexitybox}[Approach 3 Complexity]
\textbf{Time.} $O(n)$ (the $4 \times 3$ factor is a constant).
\textbf{Space.} $O(n)$ table, no recursion stack.
\end{complexitybox}

\subsection*{Approach 4: Space Optimization}

\begin{lstlisting}
#include <bits/stdc++.h>
using namespace std;

int ninjaTrainingSpaceOptimized(vector<vector<int>>& points) {
    int n = points.size();
    vector<int> prevDay(4, 0);

    for (int last = 0; last < 4; last++) {
        int best = 0;
        for (int task = 0; task < 3; task++) {
            if (task != last) best = max(best, points[0][task]);
        }
        prevDay[last] = best;
    }

    for (int day = 1; day < n; day++) {
        vector<int> currentDay(4, 0);
        for (int last = 0; last < 4; last++) {
            int best = 0;
            for (int task = 0; task < 3; task++) {
                if (task != last) {
                    best = max(best, points[day][task] + prevDay[task]);
                }
            }
            currentDay[last] = best;
        }
        prevDay = currentDay;
    }
    return prevDay[3];
}

int main() {
    vector<vector<int>> points = {{10,40,70},{20,50,80},{30,60,90}};
    cout << ninjaTrainingSpaceOptimized(points) << endl; // 210
    return 0;
}
\end{lstlisting}

\begin{complexitybox}[Approach 4 Complexity]
\textbf{Time.} $O(n)$.
\textbf{Space.} $O(1)$ (a fixed-size array of $4$ entries, independent
of $n$).
\end{complexitybox}

\begin{takeawaybox}
A DP state's second dimension does not need to be spatial -- here, it is
``which choice is currently forbidden,'' a small, fixed-size constraint
dimension exactly analogous to the visited-tracking state used in
backtracking, but folded directly into the DP table's indexing instead
of an external visited structure. Whenever a problem's choice at step
$i$ is constrained only by the choice at step $i-1$ (not by anything
further back), this ``day $\times$ last choice'' state shape applies
directly.
\end{takeawaybox}

\section{Ninja and his Friends (3D DP / Cherry Pickup)}
\label{sec:dp-ninja-friends}

\problemstatement{Given an $m \times n$ grid of chocolates, two
collectors start at the top row -- one at $(0,0)$, the other at
$(0,n-1)$ -- and both move \textbf{simultaneously}, one row at a time,
down to row $m-1$. From $(\textit{row},\textit{col})$, a collector may
move to $(\textit{row}+1,\textit{col}-1)$,
$(\textit{row}+1,\textit{col})$, or $(\textit{row}+1,\textit{col}+1)$.
Each collects the chocolates in every cell it visits; if both collectors
occupy the \textbf{same} cell at the same time, that cell's chocolates
are counted only \textbf{once}. Maximize the total chocolates collected.}

\begin{patternbox}
This closing problem of the chapter is a deliberate capstone: \emph{``two
agents moving simultaneously through the same grid''} is
Section~\ref{sec:dp-grids-intro}'s 3D DP keyword pattern. The state must
track \textbf{three} coordinates -- the shared row (since both move in
lockstep, one row per step) and \emph{each} collector's column -- because
the optimal move for one collector can depend on where the other
currently is (specifically, whether they currently coincide, which
determines whether a cell's chocolates are double-counted).
\end{patternbox}

\subsection*{Deriving the recurrence}

Let $f(\textit{row}, \textit{col1}, \textit{col2})$ be the maximum
chocolates collectible from row $\textit{row}$ down to the last row,
given that collector $1$ is currently at column $\textit{col1}$ and
collector $2$ is currently at column $\textit{col2}$ (both in the same
row, since they move together). The chocolates gained at this row are
$\textit{grid}[\textit{row}][\textit{col1}]$ plus
$\textit{grid}[\textit{row}][\textit{col2}]$ \emph{unless}
$\textit{col1}=\textit{col2}$, in which case the cell's value is added
only once. Each collector independently chooses one of $3$ column
moves ($-1, 0, +1$) for the next row, giving $3 \times 3 = 9$
combinations to maximize over:
\[
  f(\textit{row}, \textit{col1}, \textit{col2}) = \textit{collected} +
  \max_{\substack{d_1, d_2 \in \{-1,0,1\} \\ \text{both columns valid}}}
  f(\textit{row}+1, \textit{col1}+d_1, \textit{col2}+d_2),
\]
where $\textit{collected} = \textit{grid}[\textit{row}][\textit{col1}] +
(\textit{col1} \ne \textit{col2} ? \textit{grid}[\textit{row}][\textit{col2}]
: 0)$. Base case: at the last row ($\textit{row}=m-1$), there is no
further move, so $f(m-1,\textit{col1},\textit{col2}) =
\textit{collected}$ directly. The overall answer is $f(0, 0, n-1)$.

\subsection*{Approach 1: Recursion}

\begin{lstlisting}
#include <bits/stdc++.h>
using namespace std;

int solve(int row, int col1, int col2, vector<vector<int>>& grid, int m, int n) {
    if (col1 < 0 || col1 >= n || col2 < 0 || col2 >= n) return INT_MIN;

    int collected = grid[row][col1];
    if (col1 != col2) collected += grid[row][col2];

    if (row == m - 1) return collected;

    int best = INT_MIN;
    for (int d1 = -1; d1 <= 1; d1++) {
        for (int d2 = -1; d2 <= 1; d2++) {
            int next = solve(row + 1, col1 + d1, col2 + d2, grid, m, n);
            if (next != INT_MIN) best = max(best, next);
        }
    }
    return collected + best;
}

int main() {
    vector<vector<int>> grid = {{3,1,1},{2,5,1},{1,5,5},{2,1,1}};
    int m = grid.size(), n = grid[0].size();
    cout << solve(0, 0, n - 1, grid, m, n) << endl; // 24
    return 0;
}
\end{lstlisting}

\begin{complexitybox}[Approach 1 Complexity]
\textbf{Time.} $O(9^m)$ -- up to $9$ branches per call, depth $m$.
\textbf{Space.} $O(m)$ recursion stack.
\end{complexitybox}

\subsection*{Approach 2: Memoization}

\begin{lstlisting}
#include <bits/stdc++.h>
using namespace std;

int solveMemo(int row, int col1, int col2, vector<vector<int>>& grid,
              int m, int n, vector<vector<vector<int>>>& memo) {
    if (col1 < 0 || col1 >= n || col2 < 0 || col2 >= n) return INT_MIN;

    int collected = grid[row][col1];
    if (col1 != col2) collected += grid[row][col2];

    if (row == m - 1) return collected;
    if (memo[row][col1][col2] != -1) return memo[row][col1][col2];

    int best = INT_MIN;
    for (int d1 = -1; d1 <= 1; d1++) {
        for (int d2 = -1; d2 <= 1; d2++) {
            int next = solveMemo(row + 1, col1 + d1, col2 + d2, grid, m, n, memo);
            if (next != INT_MIN) best = max(best, next);
        }
    }
    return memo[row][col1][col2] = collected + best;
}

int main() {
    vector<vector<int>> grid = {{3,1,1},{2,5,1},{1,5,5},{2,1,1}};
    int m = grid.size(), n = grid[0].size();
    vector<vector<vector<int>>> memo(m, vector<vector<int>>(n, vector<int>(n, -1)));
    cout << solveMemo(0, 0, n - 1, grid, m, n, memo) << endl; // 24
    return 0;
}
\end{lstlisting}

\begin{complexitybox}[Approach 2 Complexity]
\textbf{Time.} $O(m \cdot n^2 \cdot 9)= O(m \cdot n^2)$ -- $m \cdot n^2$
distinct states, each doing $O(9)=O(1)$ work.
\textbf{Space.} $O(m \cdot n^2)$ memo $+$ $O(m)$ recursion stack.
\end{complexitybox}

\subsection*{Approach 3: Tabulation}

\begin{lstlisting}
#include <bits/stdc++.h>
using namespace std;

int solveTabulation(vector<vector<int>>& grid) {
    int m = grid.size(), n = grid[0].size();
    vector<vector<vector<int>>> dp(m, vector<vector<int>>(n, vector<int>(n, 0)));

    for (int col1 = 0; col1 < n; col1++) {
        for (int col2 = 0; col2 < n; col2++) {
            dp[m-1][col1][col2] = grid[m-1][col1] +
                (col1 != col2 ? grid[m-1][col2] : 0);
        }
    }

    for (int row = m - 2; row >= 0; row--) {
        for (int col1 = 0; col1 < n; col1++) {
            for (int col2 = 0; col2 < n; col2++) {
                int collected = grid[row][col1] +
                    (col1 != col2 ? grid[row][col2] : 0);

                int best = INT_MIN;
                for (int d1 = -1; d1 <= 1; d1++) {
                    for (int d2 = -1; d2 <= 1; d2++) {
                        int nc1 = col1 + d1, nc2 = col2 + d2;
                        if (nc1 >= 0 && nc1 < n && nc2 >= 0 && nc2 < n) {
                            best = max(best, dp[row+1][nc1][nc2]);
                        }
                    }
                }
                dp[row][col1][col2] = collected + (best == INT_MIN ? 0 : best);
            }
        }
    }
    return dp[0][0][n-1];
}

int main() {
    vector<vector<int>> grid = {{3,1,1},{2,5,1},{1,5,5},{2,1,1}};
    cout << solveTabulation(grid) << endl; // 24
    return 0;
}
\end{lstlisting}

\subsection*{Dry run}

\dryrun{Take $\textit{grid} = \begin{bmatrix}3&1&1\\2&5&1\\1&5&5\\2&1&1\end{bmatrix}$.}

The last row's table is filled directly: e.g.\
$\textit{dp}[3][0][2] = \textit{grid}[3][0]+\textit{grid}[3][2] =
2+1=3$ (collectors at different columns, both counted). Working
upward, $\textit{dp}[2][\cdot][\cdot]$ combines row $2$'s chocolates
with the best reachable entry of $\textit{dp}[3]$, and so on up to
$\textit{dp}[0][0][2]$. Tracing the full optimal path (confirmed by the
code): collector $1$ follows $0\to1\to2\to1$ (columns), collector $2$
follows $2\to1\to2\to1$, jointly visiting cells that sum to a maximum
of $\boxed{24}$ chocolates.

\begin{complexitybox}[Approach 3 Complexity]
\textbf{Time.} $O(m \cdot n^2)$.
\textbf{Space.} $O(m \cdot n^2)$ table, no recursion stack.
\end{complexitybox}

\subsection*{Approach 4: Space Optimization}

\begin{ideabox}[Only the Next Row's 2D Slice Is Ever Needed]
Exactly as in Section~\ref{sec:dp-unique-paths}, row $\textit{row}$'s
table only ever reads from row $\textit{row}+1$ -- so the full 3D table
collapses to two rolling 2D slices, each of size $n \times n$, instead of
$m \times n \times n$.
\end{ideabox}

\begin{lstlisting}
#include <bits/stdc++.h>
using namespace std;

int solveSpaceOptimized(vector<vector<int>>& grid) {
    int m = grid.size(), n = grid[0].size();
    vector<vector<int>> nextRow(n, vector<int>(n, 0));

    for (int col1 = 0; col1 < n; col1++)
        for (int col2 = 0; col2 < n; col2++)
            nextRow[col1][col2] = grid[m-1][col1] +
                (col1 != col2 ? grid[m-1][col2] : 0);

    for (int row = m - 2; row >= 0; row--) {
        vector<vector<int>> currentRow(n, vector<int>(n, 0));
        for (int col1 = 0; col1 < n; col1++) {
            for (int col2 = 0; col2 < n; col2++) {
                int collected = grid[row][col1] +
                    (col1 != col2 ? grid[row][col2] : 0);

                int best = INT_MIN;
                for (int d1 = -1; d1 <= 1; d1++) {
                    for (int d2 = -1; d2 <= 1; d2++) {
                        int nc1 = col1 + d1, nc2 = col2 + d2;
                        if (nc1 >= 0 && nc1 < n && nc2 >= 0 && nc2 < n) {
                            best = max(best, nextRow[nc1][nc2]);
                        }
                    }
                }
                currentRow[col1][col2] = collected + (best == INT_MIN ? 0 : best);
            }
        }
        nextRow = currentRow;
    }
    return nextRow[0][n-1];
}

int main() {
    vector<vector<int>> grid = {{3,1,1},{2,5,1},{1,5,5},{2,1,1}};
    cout << solveSpaceOptimized(grid) << endl; // 24
    return 0;
}
\end{lstlisting}

\begin{complexitybox}[Approach 4 Complexity]
\textbf{Time.} $O(m \cdot n^2)$.
\textbf{Space.} $O(n^2)$ for two rolling 2D slices, instead of
$O(m \cdot n^2)$ for the full 3D table.
\end{complexitybox}

\begin{takeawaybox}
A 3D DP state is exactly what ``$k$ agents moving through the same
structure simultaneously'' demands: one shared dimension for the
synchronized step (here, the row both collectors are always on
together), plus one dimension per agent for its own position. Everything
about the four-stage methodology -- recursion first, then memoization,
then tabulation in the base case's direction, then collapsing to rolling
slices -- is unchanged from the simpler 1D and 2D problems earlier in
this book; only the number of dimensions being tracked grows.
\end{takeawaybox}

